\chapter{Definition of Foundational Theories}
\label{chap:definition-of-foundational-theories}


\section{Physics vs. Philosophy}

Any argument that cannot be derived from its axioms, observed and experimentally verified, 
is classified as philosophy rather than physics.

However, any theory that attempts to explain everything cannot avoid philosopy: it must explain and predict the philosophy.

\begin{itemize}
  
\item textbf{Metaphysics (The Study of Reality)}: Explores the fundamental nature of reality, existence, and the universe.

  Key questions: What is real? What is the nature of time and space? 

\item textbf{Epistemology (The Study of Knowledge)}: Examines the nature, origin, and limits of human knowledge.

    Key questions: What is knowledge? How do we know what we know? Can we truly be certain of anything?

  \item {Axiology & Ethics (The Study of Value and Morality)}: Ethics investigates what is right and wrong, good and bad, and how we ought to live.

    Key questions: What makes an action morally right? What is justice? What is beauty?

\item {Logic (The Study of Reasoning)}: The rules of valid inference and critical thinking, providing the toolbox to construct arguments.

    Key questions: What makes an argument valid or invalid? 

\item {Philosophy of Mind}: This is the dedicated home for the study of consciousness, mental states, intentionality, and qualia (the raw, subjective "what-it-is-like" quality of experience).

  Key questions: The Hard Problem of Consciousness, coined by philosopher David Chalmers. This is the central puzzle of the field:
  How and why do physical processes in the brain give rise to subjective inner experience?
    
\end{itemize}

If the theory does not give answers to these foundational questions, then it is quite incomplete theory of everything.


\section{Minimal Requirements}

To build a fundamental theory of physics, one needs to know the minimum requirements
for a theory to qualify as a fundamental one.

Based on conventional view, such a theory must answer three basic questions:
what exists, how likely different possibilities are, and how complex structures arise from simple rules.
These three questions are often referred to as ontology, measure, and emergence, respectively.


\subsection*{1. Ontology (What exists)}

These are the basic building blocks of reality, well-defined space of possible states of the system.
For example, in Quantum Mechanics, the basic objects are quantum fields (or state vectors in Hilbert space) defined over spacetime.


\subsection*{2. Measure (What is typical or likely)}

A theory must also explain what is typical among all possible states, to allow for predictions and explanations of observed phenomena.
For example, in Quantum Field Theory, the Born rule assigns probabilities to measurement outcomes
based on the quantum state. This tells us which outcomes are likely to be observed.

\subsection*{3. Emergence (How structure arises)}

A fundamental theory must explain how complex structures arise from the basic building blocks - the ontology.
For example, chemistry emerges from Quantum Field Theory.
atoms form molecules, and molecules form complex structures, even though the underlying description
is entirely in terms of quantum fields and their interactions.


\subsection*{Hierarchy}

Without ontology, the theory has nothing to describe. Without a measure, the theory cannot make predictions.
Without emergence, the theory cannot connect to the macroscopic hierarchical structures we observe.

As with any well-designed modular software, the three layers should be mutually orthogonal without cyclic dependencies.
Ontology defines the domain. Measure provides weighting on that domain, so it depends on the ontology.
Emergence depends on both ontology and measure. 

\[
\begin{array}{c}
\mathbf{[Emergence]} \\
\downarrow \ \text{\small{Depends on both}} \\
\mathbf{[Measure]} \\
\downarrow \ \text{\small{Depends on Ontology}} \\
\mathbf{[Ontology]} \\
\end{array}
\]

Given a sufficiently rich ontology together with a suitable measure or dynamical law, higher-level structure may arise automatically.
In that sense, emergence can often be viewed as a consequence rather than an independent primitive.


\section{Classifying Current Theories}

So that's the the ideal theory. However, none of the candidate theories implement the three ingredients. 

\subsection*{Ontology (What Exists)}
\begin{itemize}
    \item \textbf{Physical primitives:} Particles, fields, spacetime (e.g., Standard Model, General Relativity)
    \item \textbf{Information-theoretic:} Binary distinctions, bits (e.g., Wheeler's ``It from Bit'')
    \item \textbf{Computational:} Programs, Turing machines or cellular automata
    \item \textbf{Mathematical:} Abstract mathematical structures (e.g., Tegmark's Mathematical Universe Hypothesis)
\end{itemize}

\subsection*{Measure (Typicality)}
\begin{itemize}
    \item \textbf{Explicit probabilistic measure:} Classical statistical mechanics (e.g., Boltzmann distributions)
    \item \textbf{Algorithmic measure:} Program-weighted probability (e.g., Solomonoff prior in Algorithmic Information Theory)
    \item \textbf{Implicit or undefined measure:} Many multiverse or mathematical universe proposals
    \item \textbf{Emergent measure:} Measure derived from structural or dynamical properties
\end{itemize}

\subsection*{Emergence}
\begin{itemize}
    \item \textbf{Dynamical laws:} Evolution governed by differential equations (e.g., GR, QFT)
    \item \textbf{Statistical emergence:} Order arising from fluctuations and typicality (e.g., Boltzmann)
    \item \textbf{Computational emergence:} Structure arising from execution of simple programs or rules (e.g., cellular automata)
    \item \textbf{Undefined or implicit emergence:} Present in many abstract frameworks without explicit mechanism
\end{itemize}


String theory / M-theory, Loop Quantum Gravity, Quantum information, Many-worlds, holographic approaches, and many more 
blurr the line between ontology and measure. It's quite a mess, really.

Looks like the architect of the universe newer took the ``Best Coding Practices''  course.



\section{Theory from First Principles}

To map this onto our earlier dilemma about reality: ontology is our attempt to define what is ``Real''. An example is Penrose's fundamental Platonic reality.
According to Penrose, mathematics is not something we invent. It exists even if we humans newer existed and discovered it. 

Conversely, measure and emergence are the ``Abstract'' mathematical machinery used to calculate predictions. This is Hawking's pragmatic toolbox.

If we accept the pragmatic path (\textbf{maximized ontology}), we accept a complex, highly specific starting state of fields, strings, or dimensions without explaining why they exist.
The reward is incredibly precise predictions.
The cost, is that the path relies on a mountain of unexplained baseline assumptions.

If we choose the ultimate path (\textbf{minimized ontology}), we gain the potential of maximal explanation.
But so far theories that belong to this group have failed give any predictions.

The price to pay is maximized abstractness.

